% lanekit report stub (vega/scripts/tests/lanekit/report_stub.py).
% main.aux and main.pdf are written directly by the stub, never compiled from this file.
\documentclass{report}
\usepackage{amsmath}
\usepackage{hyperref}
\usepackage{cleveref}
\begin{document}
\title{Stub report of lane L01}
\maketitle
\setcounter{chapter}{2}
\chapter{Report roles}
\begin{figure}\caption{intro.lane}\label{vega:intro.lane}\end{figure}
\section{chapter.forge-data}\label{vega:chapter.forge-data}
\section{chapter.lens}\label{vega:chapter.lens}
\section{chapter.shape}\label{vega:chapter.shape}
\section{chapter.blueprint}\label{vega:chapter.blueprint}
\section{chapter.net}\label{vega:chapter.net}
\section{chapter.train}\label{vega:chapter.train}
\section{chapter.eval}\label{vega:chapter.eval}
\section{chapter.pack}\label{vega:chapter.pack}
\begin{figure}\caption{shape.lineage}\label{vega:shape.lineage}\end{figure}
\begin{figure}\caption{blueprint.modules}\label{vega:blueprint.modules}\end{figure}
\begin{figure}\caption{net.blocks}\label{vega:net.blocks}\end{figure}
\begin{table}\caption{net.decisions}\label{vega:net.decisions}\end{table}
\begin{equation}\label{vega:net.forward} y = f(x) \end{equation}
\begin{equation}\label{vega:net.loss} y = f(x) \end{equation}
\begin{equation}\label{vega:net.update} y = f(x) \end{equation}
\begin{figure}\caption{train.ladder}\label{vega:train.ladder}\end{figure}
\begin{figure}\caption{train.curves}\label{vega:train.curves}\end{figure}
\begin{table}\caption{eval.metrics}\label{vega:eval.metrics}\end{table}
\begin{figure}\caption{eval.calibration}\label{vega:eval.calibration}\end{figure}
\begin{figure}\caption{pack.layout}\label{vega:pack.layout}\end{figure}
\begin{figure}\caption{pack.slice}\label{vega:pack.slice}\end{figure}
\section{conclusion}\label{vega:conclusion}
\end{document}
